\section{准粒子顺排与横/纵向分支}\label{sec:3-3}

在第3.2节中，我们建立了无未配对价核子的纯三轴刚体转子绕最大惯量轴转动的局域谐振微扰理论。然而在奇 $A$ 核以及多准粒子激发态中，原子核低能动力学不再由纯集体核芯独占，而是由未配对的价准粒子角动量 $\hat{\mathbf{j}}$ 与变形三轴核芯集体转动角动量 $\hat{\mathbf{R}}$ 发生矢量相干耦合所主导。在微观壳模型与配对平均场背景下，费米面附近的侵入高 $j$ 轨道（如稀土区 $\pi i_{13/2}$ 或 $A \approx 130$ 核区 $\pi h_{11/2}$）具有强烈的科里奥利顺排效应。价核子角动量矢量在内禀空间取向的改变，将深刻重塑转动能量曲面的极值拓扑结构与横向振动响应。

本节系统阐述奇 $A$ 三轴核中角动量矢量耦合几何与主轴惯量参考坐标的定义，构建冻结顺排（Frozen Alignment, FA）教学哈密顿量，严格展开双横向有效恢复力曲率并推导转子—准粒子耦合摆动频率公式。深入剖析纵向摆动（Longitudinal Wobbling, LW）与横向摆动（Transverse Wobbling, TW）的动力学分叉机制，对比分析 $^{163}\mathrm{Lu}$ 超形变带与 $A \approx 130$ 区正常形变候选核的实验条件差异，并确立微观重排、自洽响应与模型竞争解释的物理边界。

\subsection{角动量矢量耦合与主轴惯量参考坐标}\label{sec:3-3-1}

在粒子—转子模型（Particle-Rotor Model, PRM）框架下，总角动量由核芯与价准粒子角动量相加。本节以 $\hat{\mathbf I}$ 和 $\hat{\mathbf R}$ 表示带作用量单位的总角动量及核芯角动量；令 $\hat{\boldsymbol{\mathcal I}}=\hat{\mathbf I}/\hbar$，并以 $\hat{\mathbf j}$ 表示已经除以 $\hbar$ 的无量纲准粒子角动量。于是
\begin{equation}\label{eq:3-3-vector-coupling}
\hat{\boldsymbol{\mathcal{I}}} = \frac{\hat{\mathbf{R}}}{\hbar} + \hat{\mathbf{j}} \iff \frac{\hat{\mathbf{R}}}{\hbar} = \hat{\boldsymbol{\mathcal{I}}} - \hat{\mathbf{j}}.
\end{equation}
设三轴核芯的集体转动动能由三个主轴惯量 $\mathcal{J}_k > 0$（$k = 1, 2, 3$）决定，相应的能量量纲转动常数定义为 $A_k \equiv \frac{\hbar^2}{2\mathcal{J}_k}$。核芯哈密顿量在内禀主轴坐标系下写作：
\begin{equation}\label{eq:3-3-core-H}
\hat{H}_{\text{core}} = \sum_{k=1}^3 A_k \left( \mathcal{I}_k - j_k \right)^2.
\end{equation}

在讨论准粒子耦合系统时，必须严格区分以下两对核心概念：
\begin{enumerate}
\item \textbf{绝对总角动量 $\hat{\mathbf{I}}$ 与核芯角动量 $\hat{\mathbf{R}}$ 的区别}：实验室系守恒且具有确定良量子数的是总角动量平方 $\hat{\mathbf{I}}^2 = \hbar^2 I(I+1)$。核芯角动量 $\hat{\mathbf{R}}$ 仅是多体体系中归属于集体自由度的分量，其模长在准粒子取向改变时并不守恒；
\item \textbf{横向摆动（TW）与纵向摆动（LW）的内禀坐标定义}：在非轴对称核物理中，TW 与 LW 的区分\textbf{严格取决于准粒子顺排矢量 $\mathbf{j}$ 相对于核芯最大转动惯量轴（$\mathcal{J}_{\text{max}}$）的空间取向}，而绝不能泛称为 $\mathbf{j}$ 垂直于总角动量 $\mathbf{I}$ 或垂直于核芯 $\mathbf{R}$：
\begin{itemize}
\item \textbf{纵向摆动（Longitudinal Wobbling, LW）}：价准粒子角动量 $\mathbf{j}$ 的优势取向平行于核芯的最大转动惯量轴（即 $\mathbf{j} \parallel \text{轴 } \mathcal{J}_{\text{max}}$）；
\item \textbf{横向摆动（Transverse Wobbling, TW）}：价准粒子角动量 $\mathbf{j}$ 的优势取向垂直于核芯最大惯量轴。常见讨论中，它沿短轴或长轴，而模型最大惯量沿中间轴；这一几何轴与惯量的对应仍须在具体构型中检查。
\end{itemize}
随着自旋改变，核芯和准粒子的贡献可以重新分配，夹角也可能改变；在特定主轴冻结顺排解中，总角动量仍可沿原主轴。不能仅凭自旋增加就预设它必然连续倾斜。
\end{enumerate}

\subsection{冻结顺排哈密顿量与双横向曲率展开}\label{sec:3-3-2}

为得到透明的局域模型，采用 Frauendorf 与 D{\"o}nau（2014）讨论的冻结顺排（Frozen Alignment, FA）近似~\cite{frauendorf-2014-transverse-wobbling}。假定价准粒子的优势取向固定在主轴1，即 $\mathbf j=(j,0,0)$，$j>0$ 是无量纲顺排投影。在接近充分顺排时，可近似取它为相应高 $j$ 轨道的自旋量子数；实际投影仍由单粒子场、转动与配对共同决定。哈密顿量式~\eqref{eq:3-3-core-H} 简化为：
\begin{equation}\label{eq:3-3-FA-H}
\hat{H} = A_1 (\mathcal{I}_1 - j)^2 + A_2 \mathcal{I}_2^2 + A_3 \mathcal{I}_3^2.
\end{equation}

在大自旋极限下（总自旋 $I \gg 1$），设体系总角动量主要沿主轴1排列（$\mathcal I_1\simeq I$）。经典领先阶取 $\sum_k\mathcal I_k^2\simeq I^2$；量子Casimir为 $\sum_k\hat{\mathcal I}_k^2=I(I+1)$。体轴分量满足 $[\mathcal I_a,\mathcal I_b]=-i\epsilon_{abc}\mathcal I_c$。在同一领先阶局域近似下，
\begin{equation}\label{eq:3-3-I1-expansion}
\mathcal{I}_1 = \sqrt{I^2 - \mathcal{I}_2^2 - \mathcal{I}_3^2} \approx I - \frac{\mathcal{I}_2^2 + \mathcal{I}_3^2}{2I}.
\end{equation}
代入式~\eqref{eq:3-3-FA-H} 的第一项并展开至横向二次项：
\begin{align}
A_1 (\mathcal{I}_1 - j)^2 &\approx A_1 \left( I - j - \frac{\mathcal{I}_2^2 + \mathcal{I}_3^2}{2I} \right)^2 \nonumber \\
&\approx A_1 (I - j)^2 - 2 A_1 (I - j) \frac{\mathcal{I}_2^2 + \mathcal{I}_3^2}{2I} \nonumber \\
&= A_1 (I - j)^2 - \bar{A}_1 (\mathcal{I}_2^2 + \mathcal{I}_3^2), \label{eq:3-3-A1-expansion}
\end{align}
其中定义了随自旋演化的有效转动参数：
\begin{equation}\label{eq:3-3-A1-bar}
\bar{A}_1 \equiv A_1 \left( 1 - \frac{j}{I} \right).
\end{equation}

将展开式~\eqref{eq:3-3-A1-expansion} 代回全哈密顿量，横向动能项合并为双横向谐振子型二次型：
\begin{equation}\label{eq:3-3-transverse-H}
\hat{H} \approx A_1 (I - j)^2 + C_2 \mathcal{I}_2^2 + C_3 \mathcal{I}_3^2,
\end{equation}
其中两个横向恢复力曲率系数精确定义为：
\begin{align}
C_2 &\equiv A_2 - \bar{A}_1 = (A_2 - A_1) + A_1 \frac{j}{I}, \label{eq:3-3-C2-def} \\
C_3 &\equiv A_3 - \bar{A}_1 = (A_3 - A_1) + A_1 \frac{j}{I}. \label{eq:3-3-C3-def}
\end{align}

根据第3.2节的体轴玻色化，取 $a=(\mathcal I_2-i\mathcal I_3)/\sqrt{2I}$。原对易子为 $[\mathcal I_2,\mathcal I_3]=-i\mathcal I_1$，在同一局域领先阶中用 $\mathcal I_1\simeq I$，才有 $[a,a^\dagger]\simeq1$。对二次型进行 Bogoliubov 变换，得到该近似下的激发能量：
\begin{equation}\label{eq:3-3-wobbling-freq-FA}
\hbar \omega_{\text{wob}} = 2I \sqrt{C_2 C_3} = 2I \sqrt{\left[ (A_2 - A_1) + A_1 \frac{j}{I} \right] \left[ (A_3 - A_1) + A_1 \frac{j}{I} \right]}.
\end{equation}
本节讨论正能激发建立在局域能量极小之上的情形，其条件是两个横向二次系数均为正：
\begin{equation}\label{eq:3-3-stability-condition}
C_2 > 0 \quad \text{且} \quad C_3 > 0.
\end{equation}
仅要求根号为实数只限制 $C_2C_3\ge0$，不能单独保证所选参考态是能量极小；零曲率处也需超出此谐振近似。

\subsection{纵向分支单调上升与横向分支软化分叉机制}\label{sec:3-3-3}

式~\eqref{eq:3-3-wobbling-freq-FA} 揭示了准粒子顺排取向如何深刻改变摆动频率随自旋的动力学演化行为。我们将纵向摆动（LW）与横向摆动（TW）的能量曲率响应进行对比。

\subsubsection{纵向摆动分支（LW）：频率随自旋单调上升}

在纵向摆动中，准粒子角动量顺排轴（轴 1）恰好是核芯的最大转动惯量轴：
\begin{equation}\label{eq:3-3-LW-condition}
\mathcal{J}_1 > \mathcal{J}_2, \mathcal{J}_3 \implies A_1 < A_2, A_3.
\end{equation}
由于 $A_2 - A_1 > 0$ 且 $A_3 - A_1 > 0$，同时由于 $j < I$，有 $A_1 \frac{j}{I} > 0$。由式~\eqref{eq:3-3-C2-def} 与式~\eqref{eq:3-3-C3-def}，两个横向恢复力曲率在全自旋区间均恒为正数：
\begin{equation}\label{eq:3-3-LW-curvatures}
C_2 > A_2 - A_1 > 0, \quad C_3 > A_3 - A_1 > 0 \quad (\forall I > j).
\end{equation}
体系始终处于严格的能量极小势阱中。当总自旋 $I$ 增加时，比值 $j/I \to 0$，有效常数 $\bar{A}_1 \to A_1$ 逐渐趋近于纯转子常数。摆动频率演化行为为：
\begin{equation}\label{eq:3-3-LW-trend}
\frac{d(\hbar \omega_{\text{wob}}^{\text{LW}})}{dI} > 0.
\end{equation}
因此，在保持固定惯量与固定顺排角动量（constant-FA）的特定模型条件下，\textbf{纵向摆动（LW）的特征能量随自旋 $I$ 呈现单调增加}。这反映了在 FA 假设下随总自旋增加，围绕最大惯量主轴的旋进离心恢复力相应增强；若在微观上计入配对减弱或核芯形状演化，实际频率曲线亦可能受到动力学偏离。

\subsubsection{横向摆动分支（TW）：特征频率软化与临界分叉}

以下选取一种横向排序：准粒子顺排轴为1（例如短轴），模型最大惯量轴为2（例如中间轴），并令第三轴惯量不大于轴1：
\begin{equation}\label{eq:3-3-TW-condition}
\mathcal{J}_2 > \mathcal{J}_1 \ge \mathcal{J}_3 \implies A_2 < A_1 \le A_3.
\end{equation}
考察横向曲率 $C_2$ 的结构：
\begin{equation}\label{eq:3-3-TW-C2}
C_2 = (A_2 - A_1) + A_1 \frac{j}{I}.
\end{equation}
由于 $A_2<A_1$，无准粒子时沿轴2方向的二次系数为负；若 $A_3>A_1$，轴1是纯转子的中间惯量轴，对应欧拉鞍点。固定的顺排投影贡献正项 $A_1j/I$，可以把这个方向的局域曲率变正：
\begin{enumerate}
\item \textbf{低自旋区陀螺稳定态}：在较低自旋区间（$I$ 略大于 $j$ 时），$j/I$ 相对较大，正定耦合项压倒负曲率项使得 $C_2 > 0$。准粒子顺排成功“锚定”并稳定了体系绕轴 1 的定常旋转；
\item \textbf{自旋演化与频率软化}：在固定 FA 参数下，$\omega_{\text{wob}}^2$ 是两个关于自旋 $I$ 的线性因子的乘积。在较低自旋区，频率可能先有所上升；但随着自旋进一步增大逼近临界分叉点 $I_c$ 时，负曲率项 $(A_2 - A_1) < 0$ 逐渐显现，横向恢复力曲率 $C_2$ 减小，在逼近分叉点时呈现特征性的反常软化下降；
\item \textbf{临界自旋分叉（Bifurcation）}：当自旋达到临界值 $I_c$ 时，横向恢复力曲率降至零：
\begin{equation}\label{eq:3-3-critical-spin}
C_2(I_c) = 0 \iff 1 - \frac{j}{I_c} = \frac{A_2}{A_1} = \frac{\mathcal{J}_1}{\mathcal{J}_2} \iff I_c = \frac{j}{1 - \mathcal{J}_1 / \mathcal{J}_2}.
\end{equation}
在 $I = I_c$ 处，围绕主轴 1 的局域摆动势阱变平，简谐摆动模式发生软化（$\hbar \omega_{\text{wob}} \to 0$）。当 $I > I_c$ 时，$C_2 < 0$，绕主轴 1 的转动失去稳定性，定常转动轴从主轴 1 倾斜转入 1-2 主轴平面内，演变为平面倾斜轴转动（Planar Tilted Axis Rotation）。需要强调的是，这一分叉产生的是主轴平面内的二维倾斜旋转，属于单准粒子顺排诱导的平面转动分叉，并不构成需要三维非共面角动量三面角支撑的手征对称破缺（手征性需要非共面三维构型，见第3.5节）。
\end{enumerate}

这一趋势可直接从平方形式看出。定义 $a_2=A_2-A_1$、$a_3=A_3-A_1$，则
\begin{equation}\label{eq:3-3-frequency-squared-factors}
\frac{(\hbar\omega_{\mathrm{wob}})^2}{4}
=(a_2 I+A_1j)(a_3 I+A_1j).
\end{equation}
当 $a_2<0,a_3\ge0$ 时，一个因子趋向零，另一个保持为正，但两者乘积在整个稳定区并不必然单调。这也是能量曲线下降只能作为条件性线索的原因。

\textbf{局部转动频率定义辨析}：在理论分析中，若将局部转动角速度定义为能量对角动量的正则微商 $\hbar \omega_{\text{rot}} \equiv \frac{dE}{dI}$，基带能量 $E_0(I) \approx A_1(I - j)^2$ 给出的转动频率严格为：
\begin{equation}\label{eq:3-3-rot-freq-diff}
\hbar \omega_{\text{rot}} = 2 A_1 (I - j) = \frac{\hbar^2 (I - j)}{\mathcal{J}_1}.
\end{equation}
式中核芯角动量 $(I - j)$ 与全核角动量 $I$ 的实质差别绝不能抹去；正是这一由于准粒子顺排扣除带来的有效转动频率滞后，驱动了横向分支独特的恢复力软化。

\subsection{实验实例对比：$^{163}\mathrm{Lu}$ TSD 与 $A\approx 130$ 候选核的条件差异}\label{sec:3-3-4}

实验核谱学针对候选手征与摆动能带，在不同质量区积累了丰富的测量数据，展现出不同的输入条件与动力学行为。

\subsubsection{原型核 $^{163}\mathrm{Lu}$ 的三轴超形变横向摆动}

奇 $A$ 核 $^{163}\mathrm{Lu}$ 是横向摆动模式最经典且被实验广泛确认的原型范例（{\O}deg{\aa}rd 等人 2001，Jensen 等人 2002）~\cite{odegard-2001-evidence-wobbling-163lu,jensen-2002-second-phonon-wobbling-163lu}：
\begin{itemize}
\item \textbf{微观组态与形变解释}：有关模型把这些TSD带与 $\pi i_{13/2}$ 准粒子及三轴大形变联系起来。短轴顺排和相应形变参数属于模型描述，不能单凭质子数或轨道名称推断严格的冻结投影；
\item \textbf{惯量与动力学条件}：横向解释需要检验该组态下的最大惯量轴及准粒子取向。流体几何排序是有用参照，但实际响应还受顺排与配对影响；它不是由一幅三轴形状直接测得的三个惯量；
\item \textbf{标志性谱学证据}：观测到的TSD伴带及较强带间 $\Delta I=1$ 电四极辐射，使其成为摆动解释的重要基准；第二激发序列的非谐性还为超出单声子近似提供检验。具体寿命、强度和数据依赖在第4.4节比较，能量趋势始终与电磁资料联合使用。~\cite{odegard-2001-evidence-wobbling-163lu,jensen-2002-second-phonon-wobbling-163lu}
\end{itemize}

\subsubsection{$A \approx 130$ 区正常形变核中的争议与竞争解释}

相较于TSD案例，正常形变区的 $^{135}\mathrm{Pr}$、$^{133}\mathrm{La}$ 等候选涉及不同组态、惯量和响应条件，需要各自检验（Matta 等人2015，Biswas 等人2019）。~\cite{matta-2015-transverse-wobbling-135pr,biswas-2019-longitudinal-wobbling-133la}
\begin{itemize}
\item \textbf{组态与势面软性}：$\pi h_{11/2}$ 构型和相对较小的形变是相关讨论的输入，但同一轨道名称不能决定所有核素的惯量差或 $\gamma$ 刚度。需对选定核素和自旋窗比较具体模型；
\item \textbf{多极混合比 $\delta$ 的反演多解性}：角分布和条件角关联的非线性响应可能允许多个混合比分支，不保证总有两个，也不保证一个大于1、另一个小于1。Lv 等对 $^{135}\mathrm{Pr}$ 的联合谱学分析支持所研究连接线以M1为主，对相应E2主导摆动解释提出挑战。\cite{lv-2022-evidence-against-wobbling-135pr} 分支应由数据及共同响应条件约束；
\item \textbf{倾斜进动（TiP）与带交叉替代解释}：理论计算表明，三维转轴倾斜进动（Tilted Precession, TiP，Lawrie 2020）~\cite{lawrie-2020-tilted-precession-wobbling} 以及相邻多准粒子组态交叉同样能够在有限自旋区产生随自旋下降的能量差。因此不能将 $A \approx 130$ 区所有低激发伴带机械归因于简单的冻结顺排横向摆动。
\end{itemize}

\textbf{实验摆动能量提取公式规范}：在束谱学中提取经验摆动能量 $\hbar \omega_{\text{wob}}(I)$ 时，通常通过从摆动伴带能级中扣除相邻 yrast 带自旋能级的插值参考能量：
\begin{equation}\label{eq:3-3-Ewob-exp-def}
E_{\text{wob}}(I) \equiv E(I, \text{wob}) - \frac{E(I-1, \text{yrast}) + E(I+1, \text{yrast})}{2}.
\end{equation}
这个定义适用于所比较序列具有相应交错自旋成员、需要相邻参考带插值的情形。应列清实际带号、自旋、宇称与插值选择；其它具有共同自旋的两带当然也可以定义同自旋能差，但那是另一个参考量，不能未经说明与式~\eqref{eq:3-3-Ewob-exp-def} 混用。

\subsection{微观重排、自洽响应与模型竞争解释}\label{sec:3-3-5}

前述冻结顺排解析模型为理解横纵向摆动提供了透明的物理图像。然而在真实原子核的完整微观动力学演化中，如下关键效应将显著修正纯转子假说：
\begin{enumerate}
\item \textbf{准粒子内禀角动量的自洽解冻}：价准粒子的角动量 $\mathbf{j}$ 并非完全刚性锚定于单一主轴，而是随着核芯旋转受到离心与配对势场的动态微扰，在三维内禀空间发生连续重取向；
\item \textbf{配对减弱与频率依赖有效惯量}：随着自旋 $I$ 增加，科里奥利力逐渐拆散核内核子配对（离心反配对效应），核芯转动惯量 $\mathcal{J}_k$ 呈现出强烈的转动频率依赖性（Matsuzaki 等人 2004）~\cite{matsuzaki-2004-nuclear-inertia-and-wobbling}；
\item \textbf{频率曲线非单调性与模式判据综合性}：即便在固定参数的微观推转 RPA 或三轴准粒子转子（QTR）计算中，频率随自旋的变化曲线亦可能因能级交叉或势阱演化呈现非单调升降。因此，单纯依靠能谱能级差随自旋的升降趋势，不足以排他性地证明或否定某种摆动模式；
\item \textbf{电磁跃迁提供额外约束}：各模式对E2、M1及其干涉的预言取决于波函数、形变、磁算符和惯量组合。固定形状的局域低声子公式常有 $1/I$ 因子，但其系数也可能随自旋变化；M1大小也不能脱离组态与有效算符统一规定。第3.8节与第5.9节将把绝对强度、相对分支和混合比放入共同的前向检验。
\end{enumerate}
